% Three-layer system architecture (included from Chapter 3)
% Inter-layer spacing anchored to fit-box bottoms (not inner nodes)
\begin{tikzpicture}[
  node distance=0.85cm and 0.75cm,
  block/.style={
    draw=shublue,
    line width=0.6pt,
    fill=white,
    rounded corners=2pt,
    minimum height=0.9cm,
    minimum width=2.5cm,
    align=center,
    font=\footnotesize,
    inner sep=4pt,
  },
  wideblock/.style={
    block,
    minimum width=11.5cm,
  },
  smallblock/.style={
    block,
    minimum width=2.2cm,
    minimum height=0.8cm,
    font=\scriptsize,
  },
  layerbox/.style={
    draw=shublue!70,
    line width=0.5pt,
    rounded corners=6pt,
    inner xsep=0.55cm,
    inner ysep=0.65cm,
  },
  layertitle/.style={
    font=\bfseries\small,
    text=shublue,
    align=center,
  },
  arr/.style={
    -{Stealth[length=2.2mm, width=1.8mm]},
    line width=0.55pt,
    draw=shublue,
  },
  dashedarr/.style={
    arr,
    dashed,
  },
  note/.style={
    font=\scriptsize,
    text=shugray,
    align=center,
  },
]

% --- vertical gaps between layer boxes ---
\pgfmathsetlengthmacro{\interlayer}{2.85cm}   % box bottom to next layer top row
\pgfmathsetlengthmacro{\titleabove}{0.38cm}   % title above each fit box

% =============================================================================
% Layer 1: Scene Analysis
% =============================================================================
\node[block] (clean) {Clean\\Reference};
\node[block, right=of clean] (noise) {Noise\\Component};
\node[block, right=of noise] (mix) {Noisy\\Mixture};

\node[wideblock, below=0.95cm of noise] (metrics) {%
  Scene Metrics Extraction\\
  \scriptsize (RMS, ZCR, band SNR, local SNR, flatness)};

\node[wideblock, below=of metrics] (anass) {%
  ANASS Parameter Guidance \& Full-Dimensional Reports};

\begin{scope}[on background layer]
  \node[layerbox, fill=shulayer1, fit=(clean)(mix)(anass)] (layer1box) {};
\end{scope}

\node[layertitle, above=\titleabove of layer1box.north, anchor=south] (layer1title)
  {Layer 1: Full-Dimensional Scene Analysis};

\draw[arr] (clean.south) -- ++(0,-0.28) -| ([xshift=-0.15cm]metrics.north);
\draw[arr] (noise.south) -- (metrics.north);
\draw[arr] (mix.south) -- ++(0,-0.28) -| ([xshift=0.15cm]metrics.north);
\draw[arr] (metrics) -- (anass);

% =============================================================================
% Layer 2: Adaptive Math Denoising  (anchored below Layer 1 box)
% =============================================================================
\coordinate (layer2anchor) at ([yshift=-\interlayer]layer1box.south);

\node[block, anchor=north, minimum width=2.8cm] (router) at (layer2anchor) {Rule Router};
\node[block, right=of router, minimum width=3.6cm] (chain) {%
  Algorithm Chain\\[1pt]\scriptsize (7 modules, composable)};
\node[block, right=of chain, minimum width=2.8cm] (omlsa) {%
  OM-LSA/MCRA\\Back-end};

\begin{scope}[on background layer]
  \node[layerbox, fill=shulayer2, fit=(router)(omlsa)] (layer2box) {};
\end{scope}

\node[layertitle, above=\titleabove of layer2box.north, anchor=south] (layer2title)
  {Layer 2: Adaptive Mathematical Denoising};

\draw[arr] (anass.south) -- ++(0,-0.55) -| (router.north);
\draw[arr] (router) -- (chain);
\draw[arr] (chain) -- (omlsa);

% =============================================================================
% Layer 3: Lightweight Deep Refinement  (anchored below Layer 2 box)
% =============================================================================
\coordinate (layer3anchor) at ([yshift=-\interlayer]layer2box.south);

\node[smallblock, anchor=north] (mathout) at (layer3anchor) {Routed\\Output};
\node[block, right=of mathout, minimum width=3.8cm] (student) {%
  TinyResidualRefiner\\[1pt]\scriptsize (optional distill refine)};
\node[smallblock, right=of student] (refined) {Refined\\Output};

\node[note, below=0.12cm of student.south] (studentnote)
  {inputs: routed STFT + noisy STFT};

\begin{scope}[on background layer]
  \node[layerbox, fill=shulayer3, fit=(mathout)(refined)(studentnote)] (layer3box) {};
\end{scope}

\node[layertitle, above=\titleabove of layer3box.north, anchor=south] (layer3title)
  {Layer 3: Lightweight Deep Refinement};

\draw[arr] (omlsa.south) -- ++(0,-0.55) -| (mathout.north);
\draw[arr] (mathout) -- (student);
\draw[arr] (student) -- (refined);

% Offline teacher in the gap between Layer 3 and Layer 4
\node[smallblock, below=1.15cm of layer3box.south, draw=shugray, dashed,
      anchor=north] (teacher) {DeepFilterNet3\\(offline teacher)};

\draw[dashedarr, draw=shugray]
  (teacher.north) -- node[note, right=0.15cm] {distill train} (student.south);

% =============================================================================
% Layer 4: Evaluation & Web Platform  (anchored below teacher)
% =============================================================================
\coordinate (layer4anchor) at ([yshift=-\interlayer]teacher.south);

\node[block, anchor=north, minimum width=2.4cm] (metrics_out) at (layer4anchor) {Metrics\\Snapshot};
\node[block, right=of metrics_out, minimum width=2.5cm] (plots) {24+ Diagnostic\\Plots};
\node[block, right=of plots, minimum width=2.4cm] (audio) {A/B Audio\\Compare};
\node[block, right=of audio, minimum width=2.4cm] (abx) {ABX Blind\\Listening};

\begin{scope}[on background layer]
  \node[layerbox, fill=shulayer4, fit=(metrics_out)(abx)] (layer4box) {};
\end{scope}

\node[layertitle, above=\titleabove of layer4box.north, anchor=south] (layer4title)
  {Evaluation \& Web Platform};

\draw[arr] (refined.south) -- ++(0,-0.55) -| (plots.north);
\draw[arr] (refined.south) -- ++(0,-0.55) -| (audio.north);
\draw[arr] (plots) -- (metrics_out);
\draw[arr] (audio) -- (abx);

\end{tikzpicture}
